\subsection{Sprint 4}

Com a retomada ainda precoce das aulas presenciais, o time se reuniu para a \textit{planning} para separar as últimas \textit{tasks}. Durante o meio tempo entre a apresentação e a \textit{planning}, um dos nossos \ac{ages} III fez um \textit{rebase} no \textit{backend} do projeto para permitir que as imagens pudessem ser salvas e também para melhorar a documentação do Swagger. Por conta disso, tivemos que refatorar a integração de cursos e módulos. Anteriormente, estávamos salvando apenas a \ac{url} da imagem, mas isso estava causando erros na integração com a \ac{aws}, o que tornou a refatoração necessária.

Dividimos essa tarefa em duas etapas. Uma parte ficou responsável pelo código que exibe as imagens na tela, enquanto a outra se encarregou de enviar a imagem selecionada pelo usuário para o banco de dados. Fiquei responsável por esta última parte.

Na \textit{sprint} anterior, aprendi muito sobre \textit{mutations} e \textit{queries} no React, o que facilitou bastante a execução da minha tarefa. Assim que ambas as partes estavam prontas, nos reunimos para ajustar os detalhes necessários e garantir que a entrega ficasse o melhor possível.

Ao contrário da \textit{sprint} anterior, consegui entregar tudo que me foi proposto de forma satisfatória. Mesmo com todas as dificuldades enfrentadas devido às aulas presenciais e o estado de calamidade pública, sinto que entreguei o melhor que pude no momento. Durante a apresentação para as \textit{stakeholders}, novamente tivemos a presença apenas de uma delas, que aprovou o projeto como um todo e ficou super satisfeita com o nosso trabalho.